\section{Packing over a surface with a handle}\label{surf:section}

The comparison lemma becomes a packing construction when the base has a
handle.  A closed one-form with a nonzero period allows a change of fiber
identification across an annulus to turn the comparison inequality into
disjoint layers.  We include the boundary conditions because the fibers
will subsequently be truncated inside geometric models.

\subsection{Horizontal forms and relative primitives}

Throughout this section, let $m\geq1$ and let
\[
 \Delta=\{p\in\R_{\geq0}^{m}:\mu_\nu(p)=b_\nu\cdot p\leq c_\nu,\
                         1\leq\nu\leq N\},
 \qquad b_\nu\in\R_{\geq0}^{m},\quad c_\nu>0,
\]
be a full-dimensional compact polytope, and put
\[
 V=\{u\in\C^m:p(u)\in\Delta\},\qquad
 p_j(u)=\pi|u_j|^2,\qquad
 \omega_V=\sum_{j=1}^m\dd p_j\wedge\dd\theta_j.
\]
The faces $\mu_\nu=c_\nu$ are called outer faces.  Smooth objects on
closed sets extend to neighborhoods.  A \emph{toric horizontal one-form}
on $\Sigma\times V$ annihilates vectors tangent to $V$ and has coefficients
depending smoothly on the base point and on $p$.  We write $\dd_C$ for
differentiation in the base variables with $p$ fixed.

For such a form on an oriented surface,
\begin{equation}\label{surf:toric-volume}
 (\omega_V+\dd\Gamma)^{m+1}
  =(m+1)\omega_V^m\wedge\dd_C\Gamma.
\end{equation}
Indeed, in base coordinates write $\Gamma=A\,\dd s+B\,\dd t$.
The only possible additional term in the top exterior power is a
multiple of
$\omega_V^{m-1}\wedge\dd_pA\wedge\dd_pB\wedge\dd s\wedge\dd t$.
It vanishes because $\dd_pA,\dd_pB$ are combinations of the $\dd p_j$:
there are too many moment differentials in the displayed expression.
Thus $\dd_C\Gamma>0$ implies symplecticity.  The identity extends across
the coordinate axes by smoothness.

\begin{lemma}[Relative area primitives]\label{surf:primitive}
Let $\Sigma$ be a compact connected oriented surface with nonempty
boundary.  A smooth family $\eta_p$ of two-forms, supported in a fixed
compact subset of $\Int(\Sigma)$ and satisfying
$\int_\Sigma\eta_p=0$, admits a smooth family of one-forms $\alpha_p$,
supported in a fixed compact subset of $\Int(\Sigma)$, with
$\dd_C\alpha_p=\eta_p$.  The choice can be made by a fixed linear
operator on the two-forms.
\end{lemma}

\begin{proof}
In a coordinate disk identified with $\R^2$, write
$\eta=h(x,y)\,\dd x\wedge\dd y$, with $h$ compactly supported and of
integral zero.  Fix $\rho\in C_c^\infty(\R)$ of integral one, and set
\[
 \begin{split}
 h_0(y)&=\int_\R h(x,y)\,\dd x,\\
 A(x,y)&=\int_{-\infty}^x
             \bigl(h(v,y)-\rho(v)h_0(y)\bigr)\,\dd v,\qquad
 B(y)=\int_{-\infty}^y h_0(v)\,\dd v.
 \end{split}
\]
Then $\alpha=A\,\dd y-\rho(x)B(y)\,\dd x$ is compactly supported and
$\dd\alpha=\eta$.  For forms supported in a fixed compact set its
support can be fixed in advance.  This is a linear operator preserving
smooth parameter dependence.

Cover the given compact support by finitely many coordinate disks and
use a fixed partition of unity.  From each summand subtract its integral
times a fixed unit-integral bump in the same disk; the local construction
treats the resulting zero-integral forms.  Join these disks to one
fixed disk by finitely many chains of overlapping coordinate disks in
$\Int(\Sigma)$.  Unit-integral bumps in successive overlaps express
the difference between each original bump and a common bump as a sum
of zero-integral forms in individual disks.  Apply the local operator
to these differences.  The coefficient of the common bump is zero
because $\int_\Sigma\eta=0$.  All choices are fixed, proving the
support and parameter assertions.
\end{proof}

\begin{lemma}[Moser with outer faces]\label{surf:moser}
Let $\Sigma$ be a compact connected oriented surface with nonempty
boundary. Let $\Omega_\lambda$, $0\leq\lambda\leq1$, be a smooth family of symplectic forms on
neighborhoods of $\Sigma\times V$, and suppose
$\partial_\lambda\Omega_\lambda=\dd\alpha_\lambda$, where
$\alpha_\lambda$ is horizontal and vanishes in a fixed collar of
$\partial\Sigma\times V$.  Suppose, near each outer face and uniformly
in $\lambda$, that its standard Hamiltonian vector field satisfies
\[
       \iota_{X_{\mu_\nu}}\Omega_\lambda=-\dd\mu_\nu.
\]
Then the Moser isotopy preserves $\Sigma\times V$ and its interior,
is the identity near the base boundary, and is defined on a
neighborhood of the compact region.
\end{lemma}

\begin{proof}
Solve $\iota_{Y_\lambda}\Omega_\lambda=-\alpha_\lambda$.
Horizontality gives $\alpha_\lambda(X_{\mu_\nu})=0$, so evaluation
on $X_{\mu_\nu}$ gives $\dd\mu_\nu(Y_\lambda)=0$ near that face.
The field is zero in the base boundary collar.  Trajectories therefore
cannot cross any defining outer level in either time direction.
This applies first to interior starting points and then to boundary
points by continuity.  All inequalities are treated simultaneously,
including at nonsimple intersections of faces; no smooth-corner
assumption is needed.

The fields extend smoothly to a common neighborhood by compactness
and nondegeneracy.  Trajectories starting in the compact region stay
there and exist throughout the parameter interval.  Continuous
dependence and compactness give flows on a neighborhood as well.
For the flow $\psi_\lambda$,
\[
 \frac{\dd}{\dd\lambda}\psi_\lambda^*\Omega_\lambda
 =\psi_\lambda^*\bigl(\dd\alpha_\lambda+
                \dd\iota_{Y_\lambda}\Omega_\lambda\bigr)=0.
\]
The inverse flow has the same preservation properties, giving
preservation of the interior.
\end{proof}

\subsection{The surface packing lemma}

\begin{lemma}[Packing over a positive-genus base]\label{surf:packing}
Let $C^o$ be an oriented surface obtained from a closed connected
surface of genus at least one by removing finitely many, but at least
one, points.  Let $\Sigma_\ell\subset\Int(\Sigma_{\ell+1})$ be an
exhaustion by compact connected surfaces with smooth nonempty boundary.
Suppose that, for all sufficiently
large $\ell$, a neighborhood of $\Sigma_\ell\times V$ carries a form
\[
 \Omega_\ell=\omega_V+\dd\Gamma_\ell,\qquad
 \dd_C\Gamma_\ell=\kappa_\ell(p)>0,
\]
where $\Gamma_\ell$ is toric horizontal.  Assume that
\[
 A_\ell(p):=\int_{\Sigma_\ell}\kappa_\ell(p)
       \longrightarrow H(p)
       \quad\hbox{uniformly on }\Delta,
\]
where $H$ is continuous and positive.  The models for different
$\ell$ need not be compatible.

Let $r_1,\ldots,r_k>0$ satisfy
\[
 \{p\geq0:\textstyle\sum_jp_j\leq r_i\}\subset\Delta,
 \qquad
 \{p\geq0:\textstyle\sum_jp_j\leq r_i\}
       \cap\{\mu_\nu=c_\nu\}=\varnothing
       \quad\hbox{for every }i,\nu.
\]
Suppose the function
\[
 P(p)=H(p)-C_r(p),\qquad
 C_r(p)=\sum_{i=1}^k(r_i-\textstyle\sum_jp_j)_+,
\]
is concave, has positive Lebesgue mean on $\Delta$, and is bounded
below by a positive constant on $\Delta\setminus\tau\Delta$ for
some $0<\tau<1$.
Then, for every choice $0<r_i'<r_i$, some model contains a disjoint
symplectic packing of the closed balls $B^{2m+2}(r_i')$ in
$\Int(\Sigma_\ell\times V)$.  Each embedding is defined on a
neighborhood of its closed ball.
\end{lemma}

\begin{proof}
We carry out the construction in five steps.

\medskip\noindent
\textbf{Step 1: retain the comparison gap on a finite surface.}
By Lemma~\ref{cmp:comparison}, there are a smooth toric $K$ and a
Hamiltonian isotopy $q_\xi$, $0\leq\xi\leq1$, with $q_0=\id$,
$q_1=q$, preserving $V$, such that $(K-P)\circ q<K$.
Its generators commute with every $\mu_\nu$ near the corresponding
outer face, uniformly in time after shrinking the collars.
Write
\[
 g=\min_{v\in V}\bigl(K(q^{-1}v)-K(v)+H(p(v))-C_r(p(v))\bigr)>0
\]
and choose
$0<\varepsilon<\frac18\min(g,\min_\Delta H)$.
Take a sufficiently large $\Sigma=\Sigma_\ell$ so that
\begin{equation}\label{surf:area-error}
                  \|A_\ell-H\|_\infty<\varepsilon
\end{equation}
and $\Int(\Sigma)$ contains an embedded torus with a disk removed.
Write $\Gamma_0=\Gamma_\ell$, $\kappa=\kappa_\ell$, and
$\Omega_0=\Omega_\ell$.
We use the comparison obtained from $H$; no concavity of the finite
area function $A_\ell-C_r$ is required.

\medskip\noindent
\textbf{Step 2: place the base area on an annulus.}
Inside the handle choose annuli $A\Subset A^+\Subset\Int(\Sigma)$
around one essential circle, and a transverse annulus $B$ around a
circle intersecting it once.  There are positive coordinates
$(s,t)$ on $A^+$, with $t\in\R/\Z$, and a smooth closed one-form
$\beta$, supported in $B$, such that
\begin{equation}\label{surf:handle-form}
 \beta|_{A^+}=f(t)\,\dd t,\qquad f\geq0,\qquad
 \int_{\R/\Z}f(t)\,\dd t=1,
\end{equation}
where $f$ is positive on one open interval and zero outside its
closure.  Use the two coordinate bands on a torus, remove a disk
away from them, and take a bump form in the transverse coordinate
of $B$.  All choices can have collars inside the handle.

Choose a smooth nondecreasing $S(s)$, zero near and below the lower
end of $A$, and one near and above its upper end.  Its derivative
is supported in $\Int(A)$.  The two-form
\[
              \chi_A=S'(s)\,\dd s\wedge\beta
\]
extends by zero to $\Sigma$, is nonnegative, and has integral one.

Choose a positive area form $\kappa_{\rm sm}(p)$ agreeing with
$\kappa(p)$ in a thin base boundary collar, independent of $p$
on $A^+$, and satisfying
\begin{equation}\label{surf:small-area}
                  \sup_{p\in\Delta}\int_\Sigma\kappa_{\rm sm}(p)
                     <\varepsilon.
\end{equation}
To do this, fix a positive area form $\lambda$ on $\Sigma$ and
write $\kappa(p)=a(y,p)\lambda$.  If $\zeta$ is one near the
boundary and zero off a thin collar disjoint from $A^+$, put
$\kappa_{\rm sm}=(\zeta a+(1-\zeta)e)\lambda$.
Uniform boundedness of $a$ allows the collar integral to be made
arbitrarily small; then take $e>0$ sufficiently small.
The construction is smooth on neighborhoods and equals $e\lambda$
on $A^+$.

Set
\begin{equation}\label{surf:H-star}
 H_*(p)=\int_\Sigma(\kappa(p)-\kappa_{\rm sm}(p)),\qquad
 \kappa_1(p)=\kappa_{\rm sm}(p)+H_*(p)\chi_A.
\end{equation}
Then $H_*>0$, $\|H_*-H\|_\infty<2\varepsilon$, and
$\kappa_1>0$.  Moreover $\kappa_1-\kappa$ has zero integral
and compact interior support.  Lemma~\ref{surf:primitive} supplies
a toric horizontal $\alpha$ with
$\dd_C\alpha=\kappa_1-\kappa$ and compact interior base support.
Initially put $\widetilde\Gamma_1=\Gamma_0+\alpha$.

Normalize the primitive on a neighborhood of $\overline A$.
Choose a one-form $\gamma$ on $A^+$, independent of $p$, with
$\dd\gamma=e\lambda$; a primitive exists on the annulus.
The form
\[
 \delta(p)=\widetilde\Gamma_1-\gamma-SH_*(p)\beta
\]
is base-closed on $A^+$.  Let $c(p)$ be its period around the
annulus.  Since $\beta$ has period one, $\delta-c(p)\beta$ has
zero period, and equals $\dd_C F$ for a smooth function $F(y,p)$
on $A^+$.  Smooth parameter dependence follows by fixing a base
point and integrating the zero-period form along paths.
Choose $\zeta_A\in C_c^\infty(A^+)$ equal to one near $\overline A$,
and define globally
\[
 \Gamma_1=\widetilde\Gamma_1-c(p)\beta-\dd_C(\zeta_AF),
\]
where the last term is extended by zero.  This is horizontal and
toric, agrees with $\Gamma_0$ near $\partial\Sigma$, and satisfies
\begin{equation}\label{surf:annular-normal-form}
 \dd_C\Gamma_1=\kappa_1,\qquad
 \Gamma_1|_A=\gamma+S(s)H_*(p)\beta .
\end{equation}
The use of $\dd_C$ in the normalization retains horizontality.
All changes of the two-form are exact on the total space.
Linear interpolation from $\Gamma_0$ to $\Gamma_1$ is symplectic
by \eqref{surf:toric-volume}, because its base curvature interpolates
between $\kappa$ and $\kappa_1$.

\medskip\noindent
\textbf{Step 3: construct a positive exact path with packing layers.}
Choose smooth positive toric functions $h_i$ satisfying
\begin{equation}\label{surf:cap-smoothing}
 h_i(p)>(r_i-\textstyle\sum_jp_j)_+,\qquad
 0<\sum_i h_i-C_r<\varepsilon.
\end{equation}
For instance, use sufficiently close smooth upper approximations
$\frac12(x+\sqrt{x^2+\eta_i^2})$ to $x_+$.
Put
\[
 J=K-H_*,\qquad
 R_*(v)=K(q^{-1}v)-J(v)-\sum_i h_i(p(v)).
\]
Equations~\eqref{surf:area-error}--\eqref{surf:cap-smoothing} give
\begin{equation}\label{surf:residual-gap}
                       R_*(v)\geq g-3\varepsilon>0.
\end{equation}

The gap says that a coefficient $G(s,v)$ of $\beta$, nondecreasing in $s$,
can run from $J(v)$ to $K(q^{-1}v)$, with increments
$h_i(p(v))$ and the positive residual $R_*(v)$.
We will change fiber coordinates from the identity at the lower end
of $A$ to $q$ at the upper end. After pullback these boundary values
become $J$ and $K=H_*+J$, so both ends match $\Gamma_1$ after the
same global correction $J\beta$. Since $J$ is independent of the
base and $\beta$ is closed, this correction leaves the base curvature
unchanged and changes the total form by the exact form $\dd(J\beta)$.
The two paths below first insert the fiber change and then arrange
the increments as disjoint layers.

With the convention $\iota_{X_Q}\omega_V=-\dd Q$, write
$\partial_\xi q_\xi=X_{Q_\xi}\circ q_\xi$.
For $0\leq\lambda\leq1$, define on $A\times V$
\[
 q_s^\lambda=q_{\lambda S(s)},\qquad
 D_\lambda(s,t,u)=(s,t,q_s^\lambda u),\qquad
 G_\lambda(s,v)=S(s)H_*(q_\lambda^{-1}v).
\]
The $s$-generator is
$Q_s^\lambda=\lambda S'(s)Q_{\lambda S(s)}$, and vanishes on
the end collars. This is the Hamiltonian lifting mechanism used in
symplectic folding; compare \cite[Section~3.1, Step~3]{SchlenkFolding}.
With our sign convention, the suspension identity is
\begin{equation}\label{surf:suspension}
 D_\lambda^*\omega_V
   =\omega_V+\dd\bigl((Q_s^\lambda\circ q_s^\lambda)\,\dd s\bigr).
\end{equation}
On a pair $(\partial_s,w)$ with $w$ vertical, both added terms
equal $-\dd_u(Q_s^\lambda\circ q_s^\lambda)(w)$; on vertical
pairs the identity follows from symplecticity of $q_s^\lambda$.

Use
\begin{equation}\label{surf:first-path}
 \Omega_\lambda^A
   =\dd\gamma+
       D_\lambda^*\bigl(\omega_V+\dd(G_\lambda\beta)\bigr)
\end{equation}
on $A\times V$, retaining $\omega_V+\dd\Gamma_1$ outside.
These forms glue.  On the lower collar, $q_s^\lambda=\id$ and
$G_\lambda=0$; on the upper collar, $q_s^\lambda=q_\lambda$ and
$G_\lambda\circ q_\lambda=H_*$.  The suspension term is zero
on both collars.
Indeed, \eqref{surf:first-path} has horizontal primitive
\[
 \gamma+(Q_s^\lambda\circ q_s^\lambda)\,\dd s
             +(G_\lambda\circ q_s^\lambda)\beta
\]
after subtracting $\omega_V$, and this equals $\Gamma_1$ on the
collars.  Its change therefore extends by zero as a global
horizontal primitive.  At $\lambda=0$ it is the form from Step~2.

For any $t$-independent function $G(s,v)$,
\begin{equation}\label{surf:one-covector}
 \bigl(\dd\gamma+\omega_V+\dd(G\beta)\bigr)^{m+1}
  =(m+1)\omega_V^m\wedge
        \bigl(\dd\gamma+\partial_sG\,\dd s\wedge\beta\bigr).
\end{equation}
Every mixed term contains $\beta$, so products of two mixed terms
vanish.  Since
$\partial_sG_\lambda=S'H_*\circ q_\lambda^{-1}\geq0$,
the forms \eqref{surf:first-path} are symplectic.

Now keep $D=D_1$ fixed.  Choose successively ordered, pairwise
disjoint closed subintervals $I_1,\ldots,I_k,I_*$ in the interior
of the $s$-interval of $A$.  Choose smooth nondecreasing functions
$\rho_i,\rho_*$, each zero below its interval, one above it, and
constant near its endpoints.  Define
\begin{equation}\label{surf:G-final}
 G_{\rm f}(s,v)
      =J(v)+\sum_{i=1}^k\rho_i(s)h_i(p(v))+\rho_*(s)R_*(v).
\end{equation}
Its $s$-derivative is nonnegative.  Its lower and upper values are
$J$ and $K\circ q^{-1}=(H_*+J)\circ q^{-1}$.
On the $i$th interval it has the toric expression
\begin{equation}\label{surf:toric-layer}
 G_{\rm f}(s,v)=a_i(p(v))+\rho_i(s)h_i(p(v)),\qquad
 a_i=J+\sum_{j<i}h_j .
\end{equation}
The possibly nontoric contribution to $\partial_sG_{\rm f}$ is supported
in $I_*$; the residual term may persist above $I_*$. All preceding
packing layers are toric.

Interpolate linearly, with $0\leq\tau\leq1$, between
$G_1=SH_*\circ q^{-1}$ and $G_{\rm f}$:
$\widehat G_\tau=(1-\tau)G_1+\tau G_{\rm f}$.
On $A\times V$ use
\[
 \widehat\Omega_\tau=\dd\gamma+
                    D^*(\omega_V+\dd(\widehat G_\tau\beta)),
\]
and outside use $\omega_V+\dd(\Gamma_1+\tau J\beta)$.
At the lower end the inside primitive changes by $\tau J\beta$;
at the upper end its pullback changes by the same $\tau J\beta$.
Thus the primitives themselves glue.  The $s$-derivative of
$\widehat G_\tau$ is nonnegative, and \eqref{surf:one-covector}
proves positivity inside $A$.  Outside, the primitive is toric
and its base curvature is unchanged, since $\dd_C(J\beta)=0$.
This second exact symplectic path ends in the form $\Omega_{\rm f}$.

\medskip\noindent
\textbf{Step 4: verify the relative Moser conditions.}
Every variation primitive is horizontal and supported away from
the base boundary: its support lies in the fixed compact interior
support from Step~2 or in $A\cup\supp\beta$.
To check the outer faces, fix $\mu=\mu_\nu$.  The comparison
isotopy preserves $\mu$ and intertwines $X_\mu$ on a common
sufficiently thin collar.  This follows from
$\{Q_\xi,\mu\}=0$ there; points in a smaller collar stay in it
because $\mu$ is constant along their trajectories.
Consequently $Q_s^\lambda\circ q_s^\lambda$,
$H_*\circ q_\lambda^{-1}$, $K\circ q^{-1}$, and their indicated
pullbacks are invariant under $X_\mu$ there.  The other functions
are toric.  For every global horizontal primitive $\Gamma$ along
the paths, we therefore have
\[
 \iota_{X_\mu}\Gamma=0,\qquad
 \mathcal L_{X_\mu}\Gamma=0,\qquad
 \iota_{X_\mu}(\omega_V+\dd\Gamma)=-\dd\mu
\]
near the face.
Lemma~\ref{surf:moser} applies.  Reparametrize each stage to be
constant near its parameter endpoints to obtain a smooth
concatenated path.  There is a diffeomorphism $\psi$, defined
on a neighborhood and preserving $\Sigma\times V$ and its
interior, with $\psi^*\Omega_{\rm f}=\Omega_0$.
The final packing will return to the original model via $\psi^{-1}$.

\medskip\noindent
\textbf{Step 5: embed a closed ball in each layer.}
Use the $D$-coordinates $(s,t,v)$ and the $i$th interval.
On the open $t$-interval where $f>0$, the coordinate $T$ with
$\dd T=\beta$, normalized using \eqref{surf:handle-form}, runs
through $(0,1)$.  Write
\[
 \dd\gamma=\ell(s,T)\,\dd s\wedge\dd T,\qquad
 L(s,T)=\int_{s_i^-}^{s}\ell(\sigma,T)\,\dd\sigma ,
\]
where $s_i^-$ is the lower endpoint of $I_i$.
Then $\ell>0$, $\partial_sL>0$, and
$\dd L\wedge\dd T=\dd\gamma$.
Set
\begin{equation}\label{surf:strip-coordinate}
        x=L+G_{\rm f}-a_i=L+\rho_i h_i .
\end{equation}
At fixed $(T,v)$ this is strictly increasing in $s$.
The form becomes
\[
                \omega_V+\dd x\wedge\dd T+\dd a_i\wedge\dd T.
\]
Set $v'=\phi_{a_i}^{T}v$, using Hamiltonian time $T$.
The result is the product form $\omega_V(v')+\dd x\wedge\dd T$.
The sign follows directly from the convention: since $a_i$ is
toric, $\theta_j'=\theta_j+T\partial_{p_j}a_i$, and
\begin{equation}\label{surf:product-sign}
                 \omega_V(v')=\omega_V(v)+\dd a_i\wedge\dd T .
\end{equation}
The flow preserves the moments and is smooth across all axes.

The lower value of $x$ is zero and the upper value is
$L(s_i^+,T)+h_i(p)>h_i(p)$.  The chart therefore contains
\begin{equation}\label{surf:available-strip}
       v'\in\Int(V),\qquad 0<T<1,\qquad 0<x<h_i(p(v')).
\end{equation}
Its inverse is smooth by $\partial_sx>0$ and the implicit
function theorem, also at fiber axes.

We use an elementary consequence of Lemma~\ref{cmp:planar-disks};
related radial-area control appears in
\cite[Section~3.1, Lemma~3.2]{SchlenkHand}.
For any $A>0$ and $\delta>0$, a neighborhood of the closed planar
disk of area $A$ admits an area-preserving embedding into the
$(x,T)$-plane such that at radial area $w=\pi|z|^2$,
\begin{equation}\label{surf:thin-planar}
                   0<x<w+\delta,\qquad 0<T<1.
\end{equation}
Take a finite increasing area grid starting below $\delta/4$,
ending above $A$, and having mesh less than $\delta/4$.
At each grid value $h$, choose a smooth disk of area $h$ inside
$0<x<h+\delta/2$, $0<T<1$, with the disks strictly nested.
These can be rounded rectangles: take their heights strictly
increasing and sufficiently close to one, and horizontal
intervals strictly nested with positive left endpoints close
to zero.  The finite nesting margins persist after rounding
and the small area corrections.  Put these disks and the source
concentric disks in a sufficiently large open ambient disk and
apply Lemma~\ref{cmp:planar-disks}.  For a point of area $w$,
the next larger grid value $h$ gives
$x<h+\delta/2<w+\delta$.  The last value exceeds $A$, so the
map is defined beyond the closed disk used here.

Choose $0<\delta<r_i-r_i'$ and use \eqref{surf:thin-planar},
with $A>r_i'$, for the last complex coordinate of
$B^{2m+2}(r_i')$.  Use the first $m$ coordinates unchanged as $v'$.
On the closed ball,
\[
       w+\sum_jp_j(v')\leq r_i',
       \qquad
       x<w+\delta<r_i-\sum_jp_j(v')<h_i(p(v')).
\]
The fiber simplex is separated from the outer faces, and the planar
image is compactly contained in $0<T<1$, $x>0$.  The upper
inequality has uniform positive slack.  Thus the product embedding
lands compactly in \eqref{surf:available-strip} and extends to an
open neighborhood of the whole closed ball.  The inverse coordinate
changes give a symplectic embedding into the $i$th layer for
$\Omega_{\rm f}$.  The layers have disjoint base intervals, so
their compact images are disjoint.  Applying $\psi^{-1}$ from
Step~4 gives the required neighborhood embeddings in the original
model.
\end{proof}
